% s-335-p335-m002; original PDF 335, 336
\Aphorism{s-335-p335-m002}{382}
Wer eine That zu simuliren nöthig hat, thut in vielen Fällen am besten, sie wirklich zu begehen: so simulirt er am erfolgreichsten. Das, was die Moral in ihrer vergröbernden und gegensätzlichen Denkweise als Wahrheit und Lüge auseinanderreisst, ist im Grunde nur ein Mehr oder Minder von solchen Einzelheiten, die einen bestimmten äusseren Schein aufrechterhalten. Es giebt nur Vorspiegelung von Thatsachen, keine Thatsachen; man weiss nie, sondern \Sperr{glaubt} — sich oder Anderen. Geschieht die Beglaubigung, die Vorspiegelung von den verschiedenen Punkten aus gleichsinnig, so reden wir von objectiver Wahrheit; differiren einzelne Punkte, so reden wir von Täuschung oder Einbildung. Ein Richter weiss, dass auch, wo alle Zeugen übereinstimmen, die Wahrheit in suspenso bleibt; der Erkenntnisskritiker weiss, dass auch, wo Gesicht Gehör Getast von Einem Object berichten, über die Existenz eines „Dinges an sich“ nichts gesagt ist. Hier, wie überall, ist der Realismus eine verwerfliche Naivität; hier, wie überall, haben Gradunterschiede an Stelle von Gegensätzen zu treten. — Mancher wird gegen diese quantitative Auffassung von Wahrheit und Lüge ein paar gute Einwände auf der Zunge haben: um so besser! Als Skeptiker, der nicht mit der Wahrheit verheirathet ist, muss man ein Auge zudrücken, wenn einmal die Wahrheit mit einem Andern geht.
\par

